% !TeX root = part5.tex
\documentclass[../final.tex]{subfiles}

\begin{document}

\practice{5}{03.10.2026}{Рассеяние упругой волны на примеси.
Квантование колебаний одномерной цепочки}
\captionsetup{font=footnotesize,labelfont=bf,skip=5pt}

\begingroup
\renewcommand{\thesection}{4.\arabic{section}}

\task[4]{Рассеяние упругой волны на атоме другой массы}

\condition
Рассмотреть бесконечную одномерную цепочку атомов массы $m$
с расстоянием $a$ между соседними узлами и жёсткостью пружин $\alpha$.
В узле $n=0$ находится примесный атом массы $M$.
Решить задачу о рассеянии падающей волны $e^{ikna}$:
найти амплитуды отражения и прохождения и исследовать связь
полюса амплитуды с локализованным состоянием.

\solution

\subsection*{Уравнения движения и падающая волна}

Для обычных атомов и примеси соответственно
\begin{align*}
    m\ddot u_n&=\alpha(u_{n+1}+u_{n-1}-2u_n),\qquad n\ne0,\\
    M\ddot u_0&=\alpha(u_1+u_{-1}-2u_0).
\end{align*}
Ищем стационарное решение $u_n(t)=U_ne^{-i\omega t}$.
Частота падающей волны удовлетворяет дисперсионному соотношению
идеальной цепочки:
\[
    \omega^2(k)=\frac{4\alpha}{m}\sin^2\frac{ka}{2}.
\]
При \rassumption{$0<k<\pi/a$} волна $e^{ikna-i\omega t}$
распространяется вправо. Обозначим
\[
    U_n=
    \begin{cases}
        e^{ikna}+R e^{-ikna},& n<0,\\
        T e^{ikna},& n\ge0.
    \end{cases}
\]
Здесь $R$ и $T$ --- комплексные амплитуды отражённой и прошедшей волн.
Падающая амплитуда принята равной единице.

\subsection*{Согласование в узлах около примеси}

Уравнение обычного атома в узле $n=-1$ требует
$U_0=1+R$, а уравнение в узле $n=1$ --- $U_0=T$.
Поэтому
\[
    T=1+R.
\]
Положим $\theta=ka$ и $\gamma=M/m$. Тогда
\[
    U_1=T e^{i\theta},\qquad
    U_{-1}=e^{-i\theta}+R e^{i\theta}.
\]
Подстановка в уравнение примеси даёт
\begin{align*}
    -M\omega^2T
       &=\alpha\bigl[T e^{i\theta}+e^{-i\theta}
                         +(T-1)e^{i\theta}-2T\bigr]\\
       &=\alpha\bigl[2T(e^{i\theta}-1)
                         +e^{-i\theta}-e^{i\theta}\bigr].
\end{align*}
Следовательно,
\[
    T\bigl[M\omega^2+2\alpha(e^{i\theta}-1)\bigr]
       =2i\alpha\sin\theta.
\]

\subsection*{Амплитуды прохождения и отражения}

Используя дисперсию, получаем
\begin{align*}
    T&=\frac{2i\alpha\sin\theta}
             {4\alpha\gamma\sin^2(\theta/2)
                  +2\alpha(e^{i\theta}-1)}\\
     &=\frac{i\cos(\theta/2)}
             {\gamma\sin(\theta/2)+i e^{i\theta/2}}\\
     &=\frac{i\cos(\theta/2)}
             {(\gamma-1)\sin(\theta/2)+i\cos(\theta/2)}
       =\frac1{1-i(\gamma-1)\tan(\theta/2)},\\
    R&=T-1
       =\frac{i(\gamma-1)\tan(\theta/2)}
             {1-i(\gamma-1)\tan(\theta/2)}.
\end{align*}
При $M=m$ примесь отсутствует: $T=1$, $R=0$.

\subsection*{Коэффициенты по потоку энергии}

Среды слева и справа одинаковы, поэтому групповые скорости
падающей и прошедшей волн совпадают. Коэффициенты прохождения
и отражения по потоку равны
\begin{align*}
    \mathcal T&=|T|^2
       =\frac1{1+(\gamma-1)^2\tan^2(ka/2)},\\
    \mathcal R&=|R|^2
       =\frac{(\gamma-1)^2\tan^2(ka/2)}
             {1+(\gamma-1)^2\tan^2(ka/2)}.
\end{align*}
Отсюда \rresult{$\mathcal R+\mathcal T=1$}.
Для длинных волн $ka\ll1$
\[
    \mathcal R\simeq\frac{(\gamma-1)^2(ka)^2}{4},\qquad
    \mathcal T\longrightarrow1.
\]
При $ka\to\pi$ и $M\ne m$ получаем $\mathcal T\to0$.
Эти предельные значения понимаются как пределы изнутри фононной зоны.

\begin{rbox}[breakable=false,left=16pt,right=16pt,top=14pt,bottom=10pt]{[]}
    \centering
    \begin{tikzpicture}[x=8cm,y=3cm,
        every node/.style={text=rtext,font=\small}]
        \draw[rtext,-{Stealth[length=2mm]}] (0,0)--(1.07,0)
            node[right] {$ka/\pi$};
        \draw[rtext,-{Stealth[length=2mm]}] (0,0)--(0,1.15)
            node[above] {$\mathcal T$};
        \draw[rtext!40,dashed] (1,0)--(1,1);
        \node[below=4pt] at (0,0) {$0$};
        \node[below=4pt] at (1,0) {$1$};
        \node[left=4pt] at (0,1) {$1$};
        \draw[rtext!65,line width=1pt] (0,1)--(1,1);
        \draw[red3,line width=1.3pt,domain=0:0.995,samples=130]
            plot (\x,{1/(1+0.25*tan(90*\x)^2)});
        \draw[p3,line width=1.3pt,domain=0:0.995,samples=130]
            plot (\x,{1/(1+tan(90*\x)^2)});
        \node[above] at (0.5,1) {$M=m$};
        \node[red3,right] at (0.58,0.78) {$M=m/2$};
        \node[p3,right] at (0.12,0.58) {$M=2m$};
    \end{tikzpicture}
    \captionof{figure}{Прохождение упругой волны через массовую примесь.
    Длинные волны проходят почти без отражения;
    около границы зоны массовый дефект отражает волну.}
    \label{fig:fksv5_transmission}
\end{rbox}

\subsection*{Полюс амплитуды и локализованная мода}

Аналитически продолжим $T(k)$ на комплексные значения $k$.
Его знаменатель обращается в нуль при
\[
    \gamma\sin\frac{ka}{2}+i e^{ika/2}=0.
\]
Умножая на $2i e^{-ika/2}$, получаем
\[
    \gamma(1-e^{-ika})-2=0,\qquad
    e^{-ika}=1-\frac2\gamma.
\]
Для затухания справа от примеси выбираем
\[
    k=\frac\pi a+i\chi,\qquad \chi>0,
    \qquad e^{-ika}=-e^{\chi a}.
\]
Следовательно,
\[
    e^{\chi a}=\frac2\gamma-1=\frac{2m}{M}-1,\qquad
    \chi a=\ln\left(\frac{2m}{M}-1\right).
\]
Условие $\chi>0$ при положительных массах выполняется только для $M<m$.
Профиль локализованной моды и её частота имеют вид
\[
    U_n=U_0(-1)^n e^{-\chi a|n|},\qquad
    \omega_{\mathrm{loc}}^2
       =\frac{4\alpha}{m\gamma(2-\gamma)}
       =\frac{4\alpha m}{M(2m-M)}.
\]
Для $0<\gamma<1$ произведение $\gamma(2-\gamma)<1$, поэтому
\[
    \omega_{\mathrm{loc}}^2>\frac{4\alpha}{m}
       =\omega_{\max}^2.
\]
Полюс амплитуды воспроизводит локализованное состояние,
найденное ранее непосредственно из уравнений движения.

\answer
\[
    T=\frac1{1-i(M/m-1)\tan(ka/2)},\qquad R=T-1,
    \qquad \mathcal R+\mathcal T=1.
\]
\rresult{Локализованная мода существует при $M<m$} и лежит выше
фононной полосы идеальной цепочки.

\clearpage
\task[18]{Квантование упругой цепочки}

\condition
Построить гамильтониан одномерной упругой цепочки
в представлении нормальных координат и в приближении вторичного
квантования. Рассмотреть $N$ одинаковых атомов массы $m$,
соединённых пружинами жёсткости $\alpha$, с периодическими
граничными условиями. Обсудить роль ангармонических членов
третьего и четвёртого порядков.

\solution

\subsection*{Гамильтониан и канонические операторы}

Обозначим смещение атома в узле $r$ через $\hat q_r$,
а сопряжённый импульс --- через $\hat p_r$.
В \rassumption{гармоническом приближении}
\[
    \hat H=\sum_{r=0}^{N-1}
       \left[\frac{\hat p_r^2}{2m}
             +\frac\alpha2(\hat q_{r+1}-\hat q_r)^2\right],
    \qquad \hat q_{r+N}=\hat q_r.
\]
Операторы смещения и импульса эрмитовы и удовлетворяют соотношениям
\[
    [\hat q_r,\hat p_s]=i\hbar\delta_{rs},\qquad
    [\hat q_r,\hat q_s]=[\hat p_r,\hat p_s]=0.
\]

\subsection*{Преобразование Фурье}

Здесь используем \rassumption{безразмерное волновое число}
$k=2\pi j/N$, $j=0,\ldots,N-1$, с выбором представителей
в зоне $[-\pi,\pi)$. Физический волновой вектор равен $k/a$.
Все равенства волновых чисел, в том числе в символах Кронекера,
понимаются по модулю $2\pi$.
\begin{align*}
    \hat q_r&=\frac1{\sqrt N}\sum_k\hat Q_k e^{ikr},
        &\hat Q_k&=\frac1{\sqrt N}\sum_r\hat q_r e^{-ikr},\\
    \hat p_r&=\frac1{\sqrt N}\sum_k\hat P_k e^{-ikr},
        &\hat P_k&=\frac1{\sqrt N}\sum_r\hat p_r e^{ikr}.
\end{align*}
Основное соотношение ортогональности
\[
    \sum_{r=0}^{N-1}e^{i(k-k')r}=N\delta_{kk'}
\]
следует из суммы геометрической прогрессии.
В частности,
\[
    \sum_r e^{i(k+k')r}=N\delta_{k,-k'}.
\]
Из эрмитовости $\hat q_r$ и $\hat p_r$ получаем
\[
    \hat Q_k^\dagger=\hat Q_{-k},\qquad
    \hat P_k^\dagger=\hat P_{-k}.
\]
Например,
\[
    \hat q_r^\dagger
       =\frac1{\sqrt N}\sum_k\hat Q_k^\dagger e^{-ikr}
       =\frac1{\sqrt N}\sum_k\hat Q_{-k}^\dagger e^{ikr}
       =\hat q_r.
\]

\subsection*{Кинетическая энергия}

Подставим преобразование Фурье:
\begin{align*}
    \sum_r\frac{\hat p_r^2}{2m}
       &=\frac1{2mN}\sum_{k,k'}\hat P_k\hat P_{k'}
           \sum_r e^{-i(k+k')r}\\
       &=\frac1{2m}\sum_k\hat P_k\hat P_{-k}.
\end{align*}
В классическом лагранжиане тот же переход даёт
\[
    E_{\mathrm{kin}}=\frac m2\sum_k\dot Q_k\dot Q_{-k}.
\]

\subsection*{Потенциальная энергия}

Разность соседних смещений равна
\[
    \hat q_{r+1}-\hat q_r
       =\frac1{\sqrt N}\sum_k\hat Q_k(e^{ik}-1)e^{ikr}.
\]
Следовательно,
\begin{align*}
    \hat V&=\frac\alpha{2N}\sum_{k,k'}
       \hat Q_k\hat Q_{k'}(e^{ik}-1)(e^{ik'}-1)
       \sum_r e^{i(k+k')r}\\
       &=\frac\alpha2\sum_k\hat Q_k\hat Q_{-k}
           (e^{ik}-1)(e^{-ik}-1)\\
       &=\alpha\sum_k(1-\cos k)\hat Q_k\hat Q_{-k}.
\end{align*}
Введём частоты нормальных мод:
\[
    \omega_k^2=\frac{2\alpha}{m}(1-\cos k)
       =\frac{4\alpha}{m}\sin^2\frac k2,
    \qquad \omega_k=\omega_{-k}.
\]

\subsection*{Лагранжиан и сопряжённые импульсы}

Классический лагранжиан в нормальных координатах имеет вид
\[
    L=\frac m2\sum_k\dot Q_k\dot Q_{-k}
       -\frac m2\sum_k\omega_k^2Q_kQ_{-k}.
\]
При дифференцировании кинетической части учитываем оба члена
с индексами $k$ и $-k$:
\[
    P_k=\frac{\partial L}{\partial\dot Q_k}=m\dot Q_{-k},
    \qquad \dot Q_k=\frac{P_{-k}}m.
\]
Поэтому импульс, канонически сопряжённый $Q_k$, связан
с производной координаты противоположного волнового числа.
Это согласуется с обратным преобразованием:
\[
    p_r=m\dot q_r
       =\frac m{\sqrt N}\sum_k\dot Q_ke^{ikr}
       =\frac1{\sqrt N}\sum_k P_{-k}e^{ikr}
       =\frac1{\sqrt N}\sum_k P_ke^{-ikr}.
\]
После преобразования Лежандра и квантования
\[
    \hat H=\sum_k\left[
       \frac{\hat P_k\hat P_{-k}}{2m}
       +\frac{m\omega_k^2}{2}\hat Q_k\hat Q_{-k}\right].
\]

\subsection*{Коммутаторы нормальных координат}

Непосредственно из определений
\begin{align*}
    [\hat Q_k,\hat P_{k'}]
       &=\frac1N\sum_{r,s}e^{-ikr+ik's}[\hat q_r,\hat p_s]\\
       &=\frac{i\hbar}N\sum_{r,s}e^{-ikr+ik's}\delta_{rs}\\
       &=\frac{i\hbar}N\sum_r e^{i(k'-k)r}
         =i\hbar\delta_{kk'}.
\end{align*}
Кроме того,
\[
    [\hat Q_k,\hat Q_{k'}]=0,\qquad
    [\hat P_k,\hat P_{k'}]=0.
\]
Таким образом, выбранные преобразования сохраняют
\rresult{канонические коммутационные соотношения}.

\subsection*{Операторы рождения и уничтожения фононов}

Для $k\ne0$, когда $\omega_k>0$, определим
\begin{align*}
    \hat b_k&=\sqrt{\frac{m\omega_k}{2\hbar}}\hat Q_k
       +\frac{i}{\sqrt{2m\hbar\omega_k}}\hat P_{-k},\\
    \hat b_k^\dagger&=\sqrt{\frac{m\omega_k}{2\hbar}}\hat Q_{-k}
       -\frac{i}{\sqrt{2m\hbar\omega_k}}\hat P_k.
\end{align*}
Канонические коммутаторы дают
\[
    [\hat b_k,\hat b_{k'}^\dagger]=\delta_{kk'},\qquad
    [\hat b_k,\hat b_{k'}]=[\hat b_k^\dagger,\hat b_{k'}^\dagger]=0.
\]
Обратные формулы
\begin{align*}
    \hat Q_k&=\sqrt{\frac{\hbar}{2m\omega_k}}
          (\hat b_k+\hat b_{-k}^\dagger),\\
    \hat P_k&=-i\sqrt{\frac{m\hbar\omega_k}{2}}
          (\hat b_{-k}-\hat b_k^\dagger).
\end{align*}
Подстановка в гамильтониан диагонализует колебательную часть:
\[
    \hat H_{\mathrm{vib}}
       =\sum_{k\ne0}\hbar\omega_k
          \left(\hat b_k^\dagger\hat b_k+\frac12\right).
\]
Оператор $\hat n_k=\hat b_k^\dagger\hat b_k$ имеет собственные значения
$n_k=0,1,2,\ldots$.
\rconcept{Фонон} --- квант нормального колебания с энергией
$\hbar\omega_k$. В гармоническом приближении моды независимы.

\subsection*{Нулевая мода}

При $k=0$ все атомы смещаются одинаково, пружины не растягиваются
и $\omega_0=0$. Это свободное движение центра масс, а не осциллятор.
Если $P_{\mathrm{tot}}=\sum_r p_r$, то
\[
    P_0=\frac{P_{\mathrm{tot}}}{\sqrt N},\qquad
    \hat H_{\mathrm{cm}}=\frac{\hat P_0^2}{2m}
       =\frac{\hat P_{\mathrm{tot}}^2}{2Nm}.
\]
При фиксированном центре масс нулевую моду исключают из
колебательной суммы.

\subsection*{Роль ангармонических членов}

Разложение энергии одной связи по относительному смещению
$\Delta q_r=q_{r+1}-q_r$ имеет вид
\[
    U(\Delta q_r)=\frac\alpha2(\Delta q_r)^2
       +\frac{g_3}{3!}(\Delta q_r)^3
       +\frac{g_4}{4!}(\Delta q_r)^4+\ldots.
\]
Члены третьего и четвёртого порядков дают добавки
\[
    \hat H_3=\frac{g_3}{3!}\sum_r(\Delta\hat q_r)^3,
    \qquad \hat H_4=\frac{g_4}{4!}\sum_r(\Delta\hat q_r)^4.
\]
С $D_k=e^{ik}-1$ их выражения в нормальных координатах равны
\begin{align*}
    \hat H_3&=\frac{g_3}{3!\sqrt N}
       \sum_{k_1,k_2,k_3}\delta_{k_1+k_2+k_3,0}
         \prod_{j=1}^3(D_{k_j}\hat Q_{k_j}),\\
    \hat H_4&=\frac{g_4}{4!N}
       \sum_{k_1,k_2,k_3,k_4}\delta_{k_1+k_2+k_3+k_4,0}
         \prod_{j=1}^4(D_{k_j}\hat Q_{k_j}).
\end{align*}
После замены $Q_k$ на операторы $b_k$ и $b_k^\dagger$ эти члены
\rconcept{связывают разные фононные моды}.
Кубическая добавка описывает трёхфононные процессы, в том числе
распад одного фонона на два и обратное слияние, если это допускают
законы сохранения. Добавка четвёртого порядка содержит четырёхфононные
процессы. Квазиимпульс сохраняется с точностью до вектора обратной
решётки; при реальном рассеянии сохраняется также энергия.
Ангармонизм приводит к взаимодействию фононов и температурной
зависимости их частот и времени жизни.

\answer
\[
    \hat H=\frac{\hat P_{\mathrm{tot}}^2}{2Nm}
       +\sum_{k\ne0}\hbar\omega_k\left(\hat n_k+\frac12\right),
    \qquad \omega_k=2\sqrt{\frac\alpha m}\left|\sin\frac k2\right|.
\]
При фиксированном центре масс остаётся только фононная сумма.
Члены третьего и четвёртого порядков нарушают независимость нормальных мод
гармонической цепочки.

\endgroup
\end{document}
